This manuscript develops a proposed proof framework for the Riemann Hypothesis through
a spectral-coherence interpretation of the functional symmetry of the Riemann zeta
function $\zeta(s)$. The central object of study is the completed zeta function
$\xi(s) = \frac{1}{2}s(s-1)\pi^{-s/2}\Gamma(s/2)\zeta(s)$, which satisfies the
functional equation $\xi(s) = \xi(1-s)$ and is symmetric under reflection about
the critical line $\Re(s) = \tfrac{1}{2}$.

The manuscript proposes a spectral-coherence structure associated with $\xi$ in
which the nontrivial zeros of $\zeta$ appear as \emph{balance points} — states at
which a prescribed cancellation condition is satisfied. The framework attempts to
isolate the balance mechanism that would force all nontrivial zeros onto the
critical line $\Re(s) = \tfrac{1}{2}$, by showing that the relevant balance
condition is structurally incompatible with $\Re(s) \ne \tfrac{1}{2}$.

Specifically, the manuscript aims to:
\begin{enumerate}
  \item Define a spectral operator $H$ whose spectral properties encode the
        symmetry of $\xi$;
  \item Establish that nontrivial zeros of $\zeta$ correspond exactly to balance
        points of $H$;
  \item Prove that all balance points of $H$ lie on $\Re(s) = \tfrac{1}{2}$.
\end{enumerate}
Steps~(2) and~(3) constitute the primary open problems in the current version of
the manuscript. The framework is presented for mathematical review, with all
unresolved steps explicitly identified.

\medskip
\noindent\textbf{Keywords:} Riemann Hypothesis; zeta function; spectral theory;
reflection symmetry; critical line; coherence; operator theory.
